% Generated from validated Article Draft Result. Do not hand-edit; revise the structured draft instead.
\noindent\textbf{AgenticVBench、推論benchmarkの測定bias、Representation Forcing。5月の論文から、agent評価、serving測定、unified multimodalという三つの研究方向を短く追う。}\autocite{src-24add4f36a45,src-99d7ea892e92,src-ea2dcbeedd27}

\subsection{5月の変化を一つのsystemとして読む}

Paper Watchでは、今月の主軸記事と重なるが独立した論点を持つ3本を取り上げる。いずれも原論文は一次資料だが、ここで紹介する性能・比較・有効性は著者報告であり、独立再現済みの結論ではない。\autocite{src-24add4f36a45,src-99d7ea892e92,src-ea2dcbeedd27}


\subsection{AgenticVBench — real-world post-production評価}

AgenticVBenchはtext、image、audio、video、planning、tool useを組み合わせるreal-world post-production workflowを100 taskで評価するbenchmarkを提案する。著者報告では最良stackも約30\%をわずかに超える程度で、現実的なmultimodal workflowが依然難しいことを示すが、model ranking自体は評価構成に依存する。\autocite{src-ea2dcbeedd27}


\subsection{Measurement Bias — production inference benchmark}

この論文はproduction LLM inference benchmarkのclient-side measurementにsystemic biasが入り得ると指摘し、Python GILがTTFT／TPOT観測へ与える影響をqueueing analysisで扱う。serving性能を比較するとき、serverだけでなくmeasurement tool自身を検証対象にする必要があるという方法論的警告である。\autocite{src-24add4f36a45}


\subsection{Representation Forcing — unified multimodal}

Representation Forcingはunified multimodal modelでmodal間のrepresentation bottleneckを減らす設計を提案する。ここでは著者の性能比較を独立事実とは扱わず、text・visionなどを一つのmodelへ統合する際にintermediate representation設計が重要な研究対象になっていることを紹介する。\autocite{src-99d7ea892e92}


\subsection{Theme Synthesis}

3本に共通するのは、モデルそのものだけでなく「どう測るか」「どこで表現を分けるか」がsystem capabilityを左右する点である。agent taskの現実性、benchmark clientの測定bias、multimodal representationのbottleneckはいずれも、数字だけでは見えにくい設計条件を前景化している。\autocite{src-24add4f36a45,src-99d7ea892e92,src-ea2dcbeedd27}


\begin{claimboundary}
Claim Boundary: 原論文は研究内容を確認する一次資料だが、benchmark結果・speedup・attack success・quality比較は著者報告であり、本稿で独立再現した結果ではない。\autocite{src-24add4f36a45,src-99d7ea892e92,src-ea2dcbeedd27}
\end{claimboundary}
